\section{Universal Hashing}\label{sec:universal}

A family $\mathcal{H}$ of hash functions from a universe $U$ to $\{0, \dots, m-1\}$
is \emph{universal} if, for every pair of distinct keys $x \neq y$, a function
chosen uniformly from $\mathcal{H}$ maps them to the same bucket with probability
at most $1/m$.

The definition says nothing about any fixed function; it constrains only the
\textbf{random choice} of a function, which is why an adversary who does not know
that choice cannot force collisions.\footnote{An adversary who can observe the
table's timing may still learn enough about the chosen function to mount an
attack, as noted in~\cite{crosby2003}.}

\begin{theorem}\label{thm:chain-length}
If $h$ is drawn from a universal family and $n$ keys are inserted, the expected
length of the chain containing a given key is at most $1 + n/m$.
\end{theorem}

The bound in Theorem~\ref{thm:chain-length} depends on the load factor
$\lambda = n/m$ and on nothing else. When $\lambda$ stays below a constant, every
lookup takes expected $O(1)$ time, where the expectation is over the choice of $h$
and not over the keys.

A simple construction takes a prime $p > |U|$ and defines
$h_{a,b}(x) = ((ax + b) \bmod p) \bmod m$ with $a \in \{1, \dots, p-1\}$ and
$b \in \{0, \dots, p-1\}$; the proof that this family is universal appears in
Section~\ref{sec:proof} and uses only the fact that $\mathbb{Z}_p$ is a field.
